\documentclass[../../../clio.tex]{subfiles}
\begin{document}

\chapter{Projection Geometry}
\label{ch:chapter-04-projection-geometry}

\section{Chapter Preface}
Projection has a long technical history: perspective drawing, tomography, astronomical imaging, and computer graphics all depend on mapping richer spaces to poorer observations. The same operation both preserves invariants and destroys information.

That duality is what makes reconstruction mathematically rich. CLIO leverages projection geometry to separate what observation determines from what remains hidden, and to quantify residual ambiguity.

The chapter develops linear and nonlinear derivations, then links geometric null directions to information-theoretic uncertainty.

\section{Derivations and Formal Results}
For a linear projection \(P:\mathbb R^n\to\mathbb R^m\), every reconstruction has decomposition
\[
x=P^\dagger y + z,\qquad z\in\ker P,
\]
where \(P^\dagger\) is the Moore--Penrose pseudoinverse.

\begin{proposition}[Minimum-norm reconstruction]
\[
x^\star=P^\dagger y
\]
solves
\[
\min_x \|x\|^2\quad\text{subject to}\quad Px=y.
\]
\end{proposition}

\begin{proof}
Using Lagrangian \(\mathcal L(x,\lambda)=\|x\|^2+\lambda^\top(Px-y)\), stationarity gives \(2x+P^\top\lambda=0\). Substituting into constraint yields normal equations whose minimum-norm solution is \(x=P^\dagger y\).
\end{proof}

\begin{definition}[Differential hidden directions]
For nonlinear \(\projection\), local ambiguity directions at \(x\) are vectors in \(\ker D\projection_x\).
\end{definition}

\begin{definition}[Projection entropy]
In discrete settings, projection entropy is residual uncertainty
\[
\projectionEntropy=H(X\mid Y),
\]
with chain rule
\[
H(X,Y)=H(Y)+H(X\mid Y),\qquad H(X\mid Y)=H(X)-I(X;Y).
\]
\end{definition}

\end{document}
